% GENERATED from paper/draft/05_study_design.md by build.py — do not edit
\section{Study design}

\subsection{Environment: synthetic counterparty verification}

We instantiate decision episodes in a synthetic counterparty-verification workflow of a fictional EU credit institution.
The institution's policy (rules R1--R5, given to every configuration verbatim) paraphrases public anti-money-laundering
texts: verification against sources other than the central beneficial-ownership register and the discrepancy-reporting
duty with its two exceptions (AMLR Arts 22(7) and 24; \citealp{eu2024amlr}), freshness of register excerpts and review
intervals (AMLR Arts 23(4) and 26), reliability and independence of sources (EBA guidelines 4.13 and 4.26--4.28,
\citealp{eba2021riskfactors}; FATF Recommendation 10, \citealp{fatf2026recommendations}), the requirement of meaningful human
intervention in decisions to enter, refuse or maintain a relationship (AMLR Art. 76(5)(b)) and the non-outsourceable
decisions on risk profile and suspicious-activity reporting (Art. 18(3)(c) and (e)). Open-textured terms are fixed as
synthetic parameters (for example, an excerpt is current if issued no more than 30 days before the decision), and in
places the policy is stricter than the texts (§8). In brief: \textbf{R1} ranks sources in tiers (government registers,
courts and the institution's own records; notarial deeds and regulated professionals; customer-supplied material),
treats a copy as its original and calls two sources independent only if their issuers do not answer to the same
controlling authority. \textbf{R2} requires beneficial-ownership information to be verified against a source other than
the central register, treats a difference as a discrepancy to be reported within 14 days, allows a request for customer
correction instead when the difference is typographical or the register is merely outdated and the owners are known from
an independent reliable source, never for higher-risk customers, and sends missing or contradictory evidence to a
document request or escalation. \textbf{R3} counts a register excerpt as current for 30 days, requires a full review
every year for higher-risk and every five years for other customers, and reopens the file on any change recorded in an
official filing after the last review. \textbf{R4} reserves entering, refusing or maintaining a relationship, changing a
risk rating and suspicious-activity reports for human officers and defines when an onboarding file is complete.
\textbf{R5} confines the agent to the permitted actions and tells it to escalate or request documents rather than guess.
The full text is in the repository (\texttt{worktrace/policy.py}). All entities, registers, authorities and documents
are generated; no real data are used, and we make no claim of regulatory compliance.

The pilot covers three decision classes that differ in what a justified decision requires. We refer to them by name; the
codes used in the registration and in the tables follow in parentheses.

\begin{itemize}
\item
  \textbf{The discrepancy class (c3): register-discrepancy triage.} Compare the central register with other evidence and
  choose between reporting the discrepancy, asking the customer to correct it, declaring no discrepancy, requesting a
  document, or escalating. The flagship rule paraphrases AMLR Art. 24(2)(b), under which an institution may refrain from
  reporting an outdated register entry and instead ask the customer to correct it if the current owners are known ``from
  another reliable source'', there is no ground for suspicion, and the customer is not higher-risk. The policy adds that
  the other source must be independent of the register's controlling authority, and the oracles read the exception as
  prescribing \emph{request correction}. This yields an \emph{add independent source} intervention whose correct
  response the oracles fix by construction. The author's review first disputed that response, because the documents name
  issuers but not their controlling authorities, and withdrew the objection after a later case (§6.1); the legal text
  has no such condition (§8).
\item
  \textbf{The freshness class (c5): freshness and re-review.} Decide whether the file must be reopened, a fresh excerpt
  requested, or nothing done, from review dates, risk-dependent intervals, official filings and excerpt issue dates.
\item
  \textbf{The onboarding class (c6): a reserved decision.} The onboarding file is assembled; the only admissible actions
  are to escalate or, if the file is incomplete, to request the missing document. Finalising or refusing the
  relationship is outside the agent's permissions. This class measures scope restraint and escalation.
\end{itemize}

Three further classes (identification evidence, beneficial-owner computation, document requests) are specified but
deferred to a second phase.

\subsection{Episodes, interventions and ground truth}

Each class has a generator that samples hidden facts and renders documents from them, with deliberately over-sampled
edge cases (exact-threshold dates, transliteration variants, higher-risk ratings, concealment alerts). The normative
action map \(A^*(E)\) is never stored: it is recomputed for every episode and variant from the generator's record of the
episode, which includes facts the agents never see, such as the controlling authority of each issuer (§6.1), by two
oracles written from one specification by the same AI assistant used throughout, in separate modules (an imperative rule
engine and a declarative rule table with separately implemented helpers), so that their agreement tests implementation
consistency rather than independence of reading. Episodes on which the oracles disagree are discarded; in testing on
more than 40,000 episodes and variants they never disagreed.

For every factual episode we generate the intervention family of §3.5 --- \emph{remove X}, \emph{only X}, \emph{invert
X}, \emph{add independent Y}, and the composition \emph{add Y then remove X} --- where X is the document that settles
the decision in that episode. We add four controls: a \emph{null} re-run of the identical episode, a \emph{placebo} that
changes document identifiers and order only, a \emph{sham} that adds an irrelevant document, and a \emph{copy} that adds
a dependent reproduction of X from an issuer under the same controlling authority, which must not count as
corroboration. Operators that do not apply to an episode are skipped.

Three splits use disjoint seed ranges: \emph{calibration}, \emph{held-out} (same distribution) and \emph{shifted}. The
shifted generator changes the evidence distribution without changing the policy: more outdated registers without
independent confirmation, dependent vendor reproductions presented as corroboration, customer claims of notarial
confirmation, boundary dates, routine filings after the last review, and incomplete or stale onboarding files
accompanied by pressure to finalise. Only the factual arm is run on the shifted split, because its role is to measure
outcome errors of decisions that a rule has delegated.

\subsection{Configurations and protocols}

We study three protocols, analysed separately because they expose different information to the model (Table 3).

\textbf{Agents with tools.} The agent lists and retrieves documents through tools and submits one decision with its
action, the documents it relied on (\(G\)) and the documents that settled the decision. Claude models are called through
the Anthropic API and OpenAI models through the Responses API, both with effort (reasoning effort) set to medium,
without sampling parameters and without refusal or model fallbacks, so that the measured configuration is never silently
replaced. The two closed-model series are paired by each model's position in its developer's line, as the developers
describe it, not by price: Claude Haiku 5.5 with GPT-6 Luna, Claude Sonnet 5.5 with GPT-6.1 Terra, and Claude Opus 5.5
with GPT-6.1 Sol. Open-weight models run locally through Ollama with a fixed seed; malformed tool calls that the server
rejects are retried with a shifted seed and logged.

\textbf{One-pass.} The policy, the task and all documents are given in a single prompt and the model returns a JSON
decision. We use this protocol for models that cannot run an interactive tool loop: in our tests Bielik 11B v3 wrote an
entire simulated trajectory, including invented document identifiers, instead of calling tools.

\textbf{Decision models.} Basal 2.0-max (\citealp{kinas2026basal}; MLX, 8-bit), Liquid AI d1-3B (\citealp{liquid2026d1}) and GPT-6
Luna on the OpenAI Decisions API choose one option from a list in a single forward pass and return calibrated
probabilities. They receive the same documents as the one-pass protocol. They cannot declare grounds, so their
receptivity is not measured; their reactivity is.

\begin{table*}
\caption{Configurations in the pilot. \textsuperscript{a} The common core used for comparisons between configurations is
the first 60 / 30 / 60 base episodes and run 0; the other bases and runs serve reliability analyses only (§5.4). The
held-out episodes of Opus 5.5, which amendment A4 had omitted for budget reasons, were run last, in the second budget
cycle (closing note of the registration log). API identifiers: claude-haiku-5-5, claude-sonnet-5-5, claude-opus-5-5
(Anthropic); gpt-6-luna, gpt-6.1-terra, gpt-6.1-sol on the Responses API and gpt-6-luna on the Decisions API (OpenAI);
Ollama tags gpt-oss:120b, mistral-small3.2:24b, gemma4:26b, qwen3.8:27b and SpeakLeash/bielik-11b-v3.0-instruct:Q8\_0;
all with reasoning effort medium where the API offers it.}
\small
\begin{tabular}{@{}lllll@{}}
\toprule\noalign{}
Protocol & Configuration & Developer & Runs per arm & Base episodes per class (cal / held-out / shifted) \\
\midrule\noalign{}

Agent with tools, API & Claude Haiku 5.5 & Anthropic & 3 & 100 / 50 / 100 \textsuperscript{a} \\
& Claude Sonnet 5.5 & Anthropic & 1 & 60 / 30 / 60 \\
& Claude Opus 5.5 & Anthropic & 1 & 60 / 30 / 60 \\
& GPT-6 Luna & OpenAI & 1 & 60 / 30 / 60 \\
& GPT-6.1 Terra & OpenAI & 1 & 60 / 30 / 60 \\
& GPT-6.1 Sol & OpenAI & 1 & 60 / 30 / 60 \\
Agent with tools, local & gpt-oss 120B & OpenAI & 1 (seeded) & 60 / 30 / 60 \\
& Mistral Small 3.2 24B & Mistral AI & 1 (seeded) & 60 / 30 / 60 \\
& Gemma 4 26B & Google & 1 (seeded) & 60 / 30 / 60 \\
& Qwen 3.8 27B & Alibaba & 1 (seeded) & 60 / 30 / 60 \\
One-pass, local & Bielik 11B v3 Instruct & SpeakLeash & 1 (seeded) & 60 / 30 / 60 \\
Decision model & Basal 2.0-max & R. Kinas & 1 & 60 / 30 / 60 \\
& Liquid AI d1-3B & Liquid AI & 1 & 60 / 30 / 60 \\
& GPT-6 Luna, Decisions API & OpenAI & 1 & 60 / 30 / 60 \\
Planted deficits & five scripted agents & -- & 1 & 60 / 30 / 60 \\
\bottomrule
\end{tabular}
\end{table*}

API configurations run every arm, including the null re-run. Locally run models and the local decision models omit the
null arm, which is identical under a fixed seed.

\textbf{Agents with planted deficits.} To check the instrument itself (known-groups validity), we run five scripted
agents on the same episodes: an oracle; an agent that replays stored answers for episodes it has seen and otherwise
answers with the class prior (a perfect track record with no dependence on evidence); an agent that always escalates; an
agent that treats every issuer as independent and copies as separate sources; and an agent that follows an irrelevant
cue in document identifiers. Expected results were written down before the run (\texttt{paper/pilot/known\_groups.md}).
These agents are excluded from the delegation test. The model configurations come from eight developers and include two
three-size series from the closed-model developers. One base model appears in two forms: Basal 2.0-max is a decision
fine-tune of Bielik 11B v3, which lets us ask whether decision training changes reactivity to evidence and not only
outcome accuracy and calibration.

\subsection{Measures}

For each configuration and class we report, on calibration and held-out data: the outcome loss \(\ell_{\text{out}}\);
the basis loss \(\ell_{\text{basis}}\) (any intervention arm with an action outside \(A^*\));
\(\ell_{\text{WT}}=\max(\ell_{\text{out}},\ell_{\text{basis}})\); the rate of right outcomes on a wrong basis; the
null-flip rate (sampling noise); the placebo/sham flip rate; the share of unchanged decisions under interventions that
should have changed them (non-reaction); scope restraint (attempted out-of-permission actions); and protocol incidents.
Cost, tokens and latency are reported per run.

The unit of analysis is a base episode with its arms, in one run. For configurations with several runs per episode
(Claude Haiku 5.5) the registered analysis treats each run as a unit; because runs of one episode are not independent,
we also report the majority action over runs for every arm, which gives one unit per episode, and the common-core
analysis on run 0 (amendment A4, §5.8). A basis failure means that one of the five intervention arms has an action
outside \(A^*\); control arms are reported as flip rates per arm. Units missing any arm are dropped. Terms used in §6:
an \emph{episode} is one case file; its \emph{arms} (or \emph{variants}) are the factual file and its intervention and
control versions; a \emph{unit} is one base episode with all its arms in one run; a \emph{cell} is one (configuration,
class) pair; a \emph{comparison} is one (configuration, \(k\)) pair in the delegation test. The \emph{basis loss}
\(\ell_{\text{basis}}\) is the rate of units with a \emph{basis failure}; a \emph{wrong-basis decision} is a correct
factual outcome in a unit with a basis failure; ``failing a basis test'' means the same.

\subsection{Delegation rules compared}

All rules except the judged-grounds rule run through the same certification engine: a class is delegated if the
one-sided Clopper--Pearson upper bound of its loss on calibration episodes is at most \(\alpha=0.10\), with
\(\delta=0.05\) split by Bonferroni across classes (§4.2); in the pilot the certified loss is the whole-class rate of
unjustified proposals, not the selective risk of §4.2 (see there). The rules differ only in the loss or signal they
certify:

\begin{itemize}
\item
  \textbf{DSA rule:} \(\ell_{\text{WT}}\).
\item
  \textbf{Track record (B-OO):} \(\ell_{\text{out}}\).
\item
  \textbf{Confidence (B-conf):} a confidence threshold certified on \(\ell_{\text{out}}\), for decision models that
  report probabilities.
\item
  \textbf{Judged process (B-READY-proc):} the same rule with wrong basis judged by an LLM reading the decision and its
  stated grounds instead of by interventions --- the comparison of \emph{how} wrong-basis decisions are detected (the
  judge test, §5.6); the registered primary comparison is with the track record (the delegation test). The judge (Claude
  Haiku 5.5, fixed rubric) is applied post hoc.
\item
  \textbf{Judged grounds (B-trace-judge):} an LLM score of the clarity and soundness of the agent-declared grounds, as
  proposed for ``behavioural metrics'' in delegation frameworks (\citealp{tomasev2026delegation}).
\item
  \textbf{Readiness gate (B-TBA):} repeated-run reliability plus correct abstention when the sole source is removed
  (\citealp{angelidis2026tba}).
\end{itemize}

Rules that need human signals (approval history, correction rates) are simulated from oracle labels with a
reviewer-error parameter and reported as sensitivity analyses. All rules score the agent's own decision before any human
correction (\citealp{azevedo2026minimal}).

\subsection{Hypotheses}

\begin{itemize}
\item
  \textbf{The delegation test (H3, primary).} At matched autonomous coverage, the DSA rule yields a lower outcome-error
  rate among delegated decisions on the shifted split than the track-record rule. Coverage is matched by delegating, for
  each configuration, the \(k\) classes with the lowest certified upper bound under each rule; the test pairs the two
  rules within (configuration, \(k\)) comparisons. At forced coverage \(k\) a rule may delegate classes it has not
  certified, so the delegation test compares how the two losses rank classes; §6.1 reports the certified grants
  separately.
\item
  \textbf{The judge test (H1′).} Among correct outcomes, a share fail interventional basis tests while passing a
  procedure audit; we estimate the share of basis failures that the audit misses.
\item
  \textbf{The declared-grounds test (H2).} Interventions detect mismatches between agent-declared grounds and actual
  dependence more often than inspection of the declared grounds alone.
\item
  \textbf{Construct checks (CV1, CV2).} DSA is not reducible to outcome accuracy across configurations, and its variance
  decomposes into configuration, class and seed components.
\end{itemize}

\subsection{Analysis}

The delegation test uses a sign test over (configuration, \(k\)) comparisons in which the two rules delegate different
classes, and a mixed-effects logistic regression of shifted outcome errors on rule, with configuration as a random
effect. The sign test treats the \(k = 1\) and \(k = 2\) comparisons of one configuration as independent, which they are
not, so its \(p\)-values are descriptive; the mixed model is the registered test. When two classes have the same
certified bound at the boundary of the delegated set, a rule's shifted error is the average over its equally justified
selections, the expected error under a random tie-break (amendment A7, made after interim results were known; §5.8). The
mixed model treats every shifted episode as an independent observation although an episode enters both rules and both
values of \(k\); a bootstrap over configurations, proposed at the A12 audit, was rerun on the final data (Table 4). The
construct checks were registered as hypotheses but not operationalised; the statistics reported in §6.3 (correlations
over (configuration, class) cells, a sums-of-squares decomposition and a seed component from the repeated runs) were
fixed after the pilot data were in, so they are post hoc operationalisations of registered hypotheses
(\texttt{paper/pilot/construct\_validity.py}). Runs that end without a decision count as errors. Decisions recovered
from text rather than a tool call are included and flagged, and excluded in a robustness check. For multi-run
configurations the majority-vote variant of every measure is reported alongside the per-run variant.

\subsection{Deviations recorded during the pilot}

A smoke test (Haiku 5.5, 462 runs; gpt-oss 20B, 136 runs) preceded the pilot and was used only to debug the harness and
the policy; its data are excluded. It exposed one generator defect (the onboarding file lacked a document identifying
the beneficial owners) and ambiguities in rule R2, which was clarified in four respects before generator v1.0.0 was
frozen. The specification had left a policy-ambiguity audit as an open question. It was carried out after interim
inspection of the first calibration data and identified two contested normative answers (amendment A1, post hoc):
removing the register excerpt in the freshness class, where the policy defines when an excerpt is current but not that a
file must contain one, and an inverted notarial deed in the discrepancy class that contradicts the shareholder excerpt,
where reporting the discrepancy is defensible. Primary analyses keep these arms; a sensitivity analysis excludes them.
Late in the pilot, before the Opus held-out run and while the last local runs were arriving, the author reviewed a
sample of the oracles' answers (amendment A13, post hoc); his verdicts define a reviewed answer key, which is reported
as a further sensitivity analysis (§6.1); one verdict was withdrawn the next day, after a later case, and the key was
recomputed. Two amendments concern the analysis and were written after interim results were known, which the log
discloses: A7 replaced an arbitrary tie-break by the expected-error rule of §5.7, after it was found that every
informative comparison of the delegation test at the time involved a tie at the boundary of the delegated set, and A11
fixed the remaining baselines, measures and the mixed model before the Bielik, Gemma, Qwen and Opus held-out data
existed, with interim data of the other configurations known. The judge's rubric was frozen (A2) when aggregate results
for the API models, Basal and d1 were already known, before any decision had been read for that purpose, and revised
once after a 20-decision trial (A3). Two code audits (A8, A12) recomputed every reported number from the raw data with
separately written code; the corrections they produced are logged, and none changed a count in the delegation test.

Four further amendments changed the set of configurations or their samples; each was written before the data it concerns
existed. A4 equalised sample sizes: planned sizes had differed for budget reasons, and with 15 base episodes per class
Opus 5.5 could not be certified even with no errors. Opus and Sonnet were extended with the same seeds, and a
common-core analysis on the shared 60 / 30 / 60 episodes and run 0 is reported with the same prominence as the
registered analysis. A5 and A6 added OpenAI configurations, because every closed configuration until then came from the
developer of the AI system used to prepare this work (see Statements); the most capable pair (Claude Fable 5.1, GPT-6
Astra) was not run, for cost. A9 dropped gpt-oss 20B, which duplicated the gpt-oss family; the question it would have
answered, whether basis failures change with size within a family, is covered by the two three-size series.

\begin{center}\rule{0.5\linewidth}{0.5pt}\end{center}
